\documentclass[11pt]{article}
\usepackage[margin=1in]{geometry}
\usepackage{amsmath,amssymb}
\usepackage[hidelinks]{hyperref}
\emergencystretch=3em

\newcommand{\F}{\mathcal F}
\newcommand{\fat}{\operatorname{fat}}
\newcommand{\vc}{\operatorname{vc}}
\newcommand{\R}{\mathbb R}

\title{Literature Status for the Rudelson--Vershynin L-infinity Entropy Exponent Question}
\author{Run fa\_banach\_001, lane 19}
\date{July 18, 2026}

\begin{document}
\maketitle

\section*{Status}

This note records that a specific open entropy-exponent question in
Rudelson--Vershynin, arXiv:math/0404192, is answered in later literature by
Aiyer--Mansour--Moran--Shao--Waknine, arXiv:2605.13684.

No new theorem is claimed here. The purpose of this packet is to clear an
open-problem signal from the queue and point to the exact later result.

\section*{Source Question}

In Section 4 of arXiv:math/0404192, after Theorem~4.4, Rudelson and Vershynin
prove the empirical $L_\infty$ entropy estimate
\[
  D_\infty(\F,t)
  \le C v \log(n/vt)\,\log^\varepsilon(2n/v),
  \qquad v=\vc(\F,c\varepsilon t).
\]
They compare this with the Alon et al. bound, which gave
$D_\infty(\F,t)=O(\log^2 n)$ in the relevant finite-domain regime, and then
state that it remains open whether the logarithmic exponent can be made $1$.

The local source evidence is PDF page~21 of \texttt{source\_paper.pdf}; in the
parsed arXiv TeX source, this sentence occurs at
\texttt{data/parsed/arxiv\_sources/0404192/source.tex}, line~1459.

\section*{Later Answer}

The later paper arXiv:2605.13684 explicitly identifies the same
Rudelson--Vershynin question. Its abstract says that the paper obtains sharp
empirical $\ell_\infty$ metric-entropy bounds and that these results resolve
open questions by Alon et al. and Rudelson--Vershynin.

The theorem-level statement is Theorem~2. For a bounded real-valued function
class $\F\subset [-R,R]^X$, let
\[
  \gamma^\ast=\inf\{\gamma>0:\fat_\gamma(\F)<\infty\}.
\]
Then the empirical metric entropy $H(\F,\gamma,n)$ is $O(\log n)$ for
$\gamma\in(2\gamma^\ast,R)$. The theorem also records the lower regimes,
including a sometimes tight $O(\log^2 n)$ behavior below the critical scale.

Immediately after Theorem~2, the authors say that, in view of the
$O(\log^{1+\varepsilon}n)$ bounds of Rudelson and Vershynin, their result
resolves the question of whether the exponent can be improved to $1$.
Therefore the source question is already answered by later literature, and the
supporting authors explicitly knew they were answering it.

\section*{Scope}

This packet only clears the $L_\infty$ empirical entropy exponent-one question
following Theorem~4.4 of arXiv:math/0404192. It does not claim to settle other
open-problem signals or neighboring geometric questions in the source paper.

The later answer is scale-sensitive: exponent-one growth holds above the
threshold in Theorem~2, while a near-$\log^2 n$ lower bound remains in the
lower regime. This sharp threshold behavior is part of the resolution.

\section*{References}

\begin{enumerate}
\item M. Rudelson and R. Vershynin,
\emph{Combinatorics of random processes and sections of convex bodies},
arXiv:math/0404192; Annals of Mathematics 164 (2006), 603--648.
\item S. Aiyer, Y. Mansour, S. Moran, H. Shao, and T. Waknine,
\emph{Scale-Sensitive Shattering: Learnability and Evaluability at Optimal
Scale}, arXiv:2605.13684.
\end{enumerate}

\end{document}
